\section{Conclusion}
\label{sec:conclusion}
We asked whether network-aware graph learning improves on independent road-segment
forecasts. On METR-LA the answer is yes, and we can say why and where. A strong
per-segment XGBoost baseline leaves errors that are correlated across adjacent
segments and predictable from the neighbours' state at forecast time
(Eq.~\eqref{eq:gap}, Fig.~\ref{fig:gap}). A graph attention network over the fixed
road graph uses exactly this information and lowers the MAE by 5\,\% overall and by
6\,\% in peak periods at 60\,min; every difference to the independent models is
significant. Treating sensor outages as missing rather than as 0\,mph is the single
most valuable design choice (0.26\,mph at 60\,min in the ablation).

The hybrid shows that a classical and a deep model can be combined soundly: the
graph network corrects the per-segment forecasts, cross-fitting prevents the
second stage from learning on optimistic in-sample residuals, and a zero-initialised
correction starts exactly at the XGBoost forecast. It is the most accurate model at
15\,min, but at longer horizons the stand-alone ST-GAT is
better, so the simple residual design does not dominate. The ablations explain where
the gains come from: outage masking and the road graph matter most (removing either
costs 0.12--0.15\,mph), the road prior beats unrestricted attention, and attention
adds little over averaging the same neighbours. In the hybrid, the graph and
cross-fitting are both essential, and under an equal, short training budget the
hybrid is more accurate than ST-GAT at every horizon.

\textbf{Limitations.} Each configuration was trained with one seed on one laptop, so
the tests and intervals capture uncertainty over test days but not over training
runs, and the ablations use a shorter budget than the main runs. METR-LA provides
speed only (no flow, occupancy, incidents or weather) and four months of data. The
attention coefficients are shared within a one-hour window and use the original,
``static'' GAT scoring. The signal experiment measures decision timing in an offline
replay, not traffic outcomes.

\textbf{Future work.} Multiple seeds; GATv2 or per-step attention; flow and
occupancy from raw PeMS station data; a learned gate between the stand-alone and the
hybrid forecast by horizon; and a calibrated SUMO or CityFlow micro-simulation to
measure what earlier signal-plan switching is worth in delay.
